% GENERATED FILE. Edit canonical lessons, then run npm run generate:books.
% canonical-source-hash: fnv1a64:d715b6eaeb0c6d50
% canonical-lessons: MR-C267-jvalamukhi, MR-C267-pathar, MR-C267-gota, MR-C267-cikhal, MR-C267-dhul

\chapter{Volcano, Plateau, Pebble, Mud, Dust}
\label{ch:mr-volcano-plateau-pebble-mud-dust}
\hlchaptermodality{267}

\begin{chapteropening}
\textbf{By the end of this chapter:} \emph{I can say a volcano, a plateau, a pebble, mud and dust.}

Complete the last lesson of chapter 267: I can say a volcano, a plateau, a pebble, mud and dust.
\end{chapteropening}

\section[\texorpdfstring{\mr{ज्वालामुखी} (jvālāmukhī)}{ज्वालामुखी (jvālāmukhī)}]{\mr{ज्वालामुखी} (jvālāmukhī) — a volcano}

\label{lesson:MR-C267-jvalamukhi}



\begin{quote}
Before the new one: say the Marathi for a licence, then the Marathi for a helmet.
\end{quote}

\subsection*{You'll want to know: \mr{ज्वालामुखी}}
\textbf{\mr{ज्वालामुखी}} — \emph{jvālāmukhī} — \textquotedblleft{}a volcano\textquotedblright{}.

\begin{grammarlens}[title={What you've built}]
One more word for this chapter.
\end{grammarlens}

\subsection*{Your turn}
\begin{itemize}
\raggedright
  \item \emph{Say it:} \emph{jvālāmukhī}
  \item \emph{Say it:} \emph{jvālāmukhī}, once more
  \item \emph{Say it:} say \emph{jvālāmukhī}
  \item \emph{Recall:} say the Marathi for a licence, then the Marathi for a helmet, then say \emph{jvālāmukhī} again
\end{itemize}

\begin{tcolorbox}[breakable,colback=teal!4,colframe=teal!35!black,title={Before you move on}]
What does \mr{ज्वालामुखी} mean? (\textquotedblleft{}A volcano\textquotedblright{}.) Say it once more.
\end{tcolorbox}
\section[\texorpdfstring{\mr{पठार} (paṭhār)}{पठार (paṭhār)}]{\mr{पठार} (paṭhār) — a plateau}

\label{lesson:MR-C267-pathar}



\begin{quote}
Before the new one: say the Marathi for a helmet, then the Marathi for a volcano.
\end{quote}

\subsection*{You'll want to know: \mr{पठार}}
\textbf{\mr{पठार}} — \emph{paṭhār} — \textquotedblleft{}a plateau\textquotedblright{}.

\begin{grammarlens}[title={What you've built}]
One more word for this chapter.
\end{grammarlens}

\subsection*{Your turn}
\begin{itemize}
\raggedright
  \item \emph{Say it:} \emph{paṭhār}
  \item \emph{Say it:} \emph{paṭhār}, once more
  \item \emph{Say it:} say \emph{paṭhār}
  \item \emph{Recall:} say the Marathi for a helmet, then the Marathi for a volcano, then say \emph{paṭhār} again
\end{itemize}

\begin{tcolorbox}[breakable,colback=teal!4,colframe=teal!35!black,title={Before you move on}]
What does \mr{पठार} mean? (\textquotedblleft{}A plateau\textquotedblright{}.) Say it once more.
\end{tcolorbox}
\section[\texorpdfstring{\mr{गोटा} (goṭā)}{गोटा (goṭā)}]{\mr{गोटा} (goṭā) — a pebble}

\label{lesson:MR-C267-gota}



\begin{quote}
Before the new one: say the Marathi for a volcano, then the Marathi for a plateau.
\end{quote}

\subsection*{You'll want to know: \mr{गोटा}}
\textbf{\mr{गोटा}} — \emph{goṭā} — \textquotedblleft{}a pebble\textquotedblright{}.

\begin{grammarlens}[title={What you've built}]
One more word for this chapter.
\end{grammarlens}

\subsection*{Your turn}
\begin{itemize}
\raggedright
  \item \emph{Say it:} \emph{goṭā}
  \item \emph{Say it:} \emph{goṭā}, once more
  \item \emph{Say it:} say \emph{goṭā}
  \item \emph{Recall:} say the Marathi for a volcano, then the Marathi for a plateau, then say \emph{goṭā} again
\end{itemize}

\begin{tcolorbox}[breakable,colback=teal!4,colframe=teal!35!black,title={Before you move on}]
What does \mr{गोटा} mean? (\textquotedblleft{}A pebble\textquotedblright{}.) Say it once more.
\end{tcolorbox}
\section[\texorpdfstring{\mr{चिखल} (cikhal)}{चिखल (cikhal)}]{\mr{चिखल} (cikhal) — mud}

\label{lesson:MR-C267-cikhal}



\begin{quote}
Before the new one: say the Marathi for a plateau, then the Marathi for a pebble.
\end{quote}

\subsection*{You'll want to know: \mr{चिखल}}
\textbf{\mr{चिखल}} — \emph{cikhal} — \textquotedblleft{}mud\textquotedblright{}.

\begin{grammarlens}[title={What you've built}]
One more word for this chapter.
\end{grammarlens}

\subsection*{Your turn}
\begin{itemize}
\raggedright
  \item \emph{Say it:} \emph{cikhal}
  \item \emph{Say it:} \emph{cikhal}, once more
  \item \emph{Say it:} say \emph{cikhal}
  \item \emph{Recall:} say the Marathi for a plateau, then the Marathi for a pebble, then say \emph{cikhal} again
\end{itemize}

\begin{tcolorbox}[breakable,colback=teal!4,colframe=teal!35!black,title={Before you move on}]
What does \mr{चिखल} mean? (\textquotedblleft{}Mud\textquotedblright{}.) Say it once more.
\end{tcolorbox}
\section[\texorpdfstring{\mr{धूळ} (dhūḷ)}{धूळ (dhūḷ)}]{\mr{धूळ} (dhūḷ) — dust}

\label{lesson:MR-C267-dhul}



\begin{quote}
Before the new one: say the Marathi for a pebble, then the Marathi for mud.
\end{quote}

\subsection*{You'll want to know: \mr{धूळ}}
\textbf{\mr{धूळ}} — \emph{dhūḷ} — \textquotedblleft{}dust\textquotedblright{}.

\begin{grammarlens}[title={What you've built}]
One more word for this chapter.
\end{grammarlens}

\subsection*{Your turn}
\begin{itemize}
\raggedright
  \item \emph{Say it:} \emph{dhūḷ}
  \item \emph{Say it:} \emph{dhūḷ}, once more
  \item \emph{Say it:} say \emph{dhūḷ}
  \item \emph{Recall:} say the Marathi for a pebble, then the Marathi for mud, then say \emph{dhūḷ} again
\end{itemize}

\begin{tcolorbox}[breakable,colback=teal!4,colframe=teal!35!black,title={Before you move on}]
What does \mr{धूळ} mean? (\textquotedblleft{}Dust\textquotedblright{}.) Say it once more.
\end{tcolorbox}
